\documentclass[10pt]{beamer-control}
\usepackage{beamer-control-singlefile}
\INCLUDEONLY{Basic control functions}
\begin{document}
\CONCEPT{Basic control functions}

\begin{SUMMARY}
\begin{itemize}
\item PID controller structure
\item P, I and D
\end{itemize}
\vfill References:
\begin{itemize}
\item \astrom{§11.1}
\item \minicite{odwyer2009-book}
\end{itemize}
\end{SUMMARY}


\SUBCONCEPT{Module \& topic introduction}

\begin{frame}
\frametitle{Control system design}
In module I, we:
\begin{itemize}
\item introduced dynamical systems and the concept of feedback
\item described differential equations using state space methods
\item represented systems in terms of inputs \& outputs, and standardised forms
\end{itemize}
In module II, we:
\begin{itemize}
\item formalised the ideas of control (reachability, etc.) 
\item introduced the frequency response (poles, zeros, Bode plot)
\item discussed the concept of stability (loop transfer function, Nyquist plot)
\end{itemize}
We can now (finally) properly discuss control system design, starting with \alert{PID Control}
\end{frame}

	
\SUBCONCEPT{PID controller structure}

\begin{frame}
\begin{figure}
	\centering
	\includegraphics[width=0.8\linewidth]{figure11.1a}
	\\
	\textbf{Figure 11.1 (a):} Block diagram of PID controller in closed loop system. 
\end{figure}
\end{frame}

\begin{frame}{Control without a model}
The PID controller is composed of three terms:
\begin{itemize}
\item {\Huge\kramer{P}}roportional term -- depends on the present error,
\item {\Huge\kramer{I}}ntegral term  -- depends on the past error, and 
\item {\Huge\kramer{D}}erivative term -- depends on the future error.
\end{itemize}
\end{frame}

\begin{frame}
\begin{figure}
	\centering
	\includegraphics{figure1.15}
	\\
	\textbf{Figure 1.15}: Action of a PID controller.
\end{figure}
\end{frame}

\begin{frame}{Control without a model}
\begin{itemize}
\item In contrast to a state feedback controller\dots
\item \dots the PID controller uses linear extrapolation of the measured output for control.
\item PID controllers are ubiquitous in industrial applications due to:
\begin{itemize}
\item their wide applicability
\item their understandability and ease of implementation
\item their avoidance of using a model [*]
\end{itemize}
\end{itemize}
\end{frame}


\begin{frame}{Why PID?}{O'Dwyer (2009) \emph{Handbook of PI and PID Controller Tuning Rules}}
PID is a standardised control methodology, with well-defined and `understandable' interfaces.
\begin{itemize}
\item Most industrial controllers are PID:
\begin{itemize}
\item For controlling low- to medium-order processes
\item When time delays are small
\item Only 5\%--10\% of plants cannot be controlled with PID
\item Large industrial plants may contain thousands of PID loops
\end{itemize}
\item But:
\begin{itemize}
\item 80\% of PID loops in industry are badly tuned
\item 25\% have `factory default' settings --- i.e., no tuning!
\item This course hopefully a small step to help this problem
\end{itemize}
\end{itemize}
\end{frame}


\begin{frame}{Control action}
\begin{itemize}
	\item The control action of a PID controller is given by
	\[u=\kkp e+\kki \int_0^t e(\tau) \, \mathrm{d} \tau+\kkd \frac{\mathrm{d} e}{\mathrm{d} t}\]
	where $\kkp$ is the proportional gain, $\kki$ is the integral gain, and $\kkd$ is the derivative gain
	\item The integral time constant $\TTi=\tfrac{\kkp}{\kki}$ and derivative time constant $\TTd = \tfrac{\kkd}{\kkp}$ provide an alternate parameterisation for the control action
	\[u = \kkp\left(e+\frac{1}{\TTi} \int_0^t e(\tau) \, \mathrm{d} \tau+\TTd \frac{\mathrm{d} e}{\mathrm{d} t}\right)\]
\end{itemize}

\end{frame}



\SUBCONCEPT{P, I and D}

\begin{frame}{Effects of the three terms}{Proportional feedback}
\begin{itemize}
\item With only P feedback, non-zero steady-state error is common
\item Increasing the proportional gain $\kkp$ will reduce the error but make the system more oscillatory
\item Too high of a gain may make the system unstable 
\item Proportional band also used, $\mathrm{PB}=\tfrac{100}{\kkp}$
\end{itemize}
\end{frame}

\begin{frame}{Effects of the three terms}{Integral feedback}
\begin{itemize}
	\item We saw in the Topic 4 `State Feedback' that integral action guarantees that if a steady state exists, we will reach it with integral control
	\item An increase in the integral gain $\kki$ makes the system robust to disturbances
	\item Too high of a gain results in oscillatory behaviour, poor robustness, and possible instability
\end{itemize}
\end{frame}


\begin{frame}{Effects of the three terms}{Derivative feedback}
	\begin{itemize}
		\item Derivative control provides predictive or anticipatory action
		\item An increase in the derivative gain $\kkd$ adds damping into the system
		\item Reduces overshoot
	\end{itemize}
\end{frame}

\begin{frame}{PID control implemented}
\begin{figure}
	\centering
	\includegraphics[width=\linewidth]{figure11.2}
	\\
	\textbf{Figure 11.2:} Responses to step changes in the reference value for a system with a proportional controller (a), PI controller (b), and PID controller (c).
\end{figure}
\end{frame}


\SUBCONCEPT{Additional details}

\begin{frame}
\frametitle{PID gains vs time constants}

\begin{itemize}
\item As illustrated, this is the parallel form of PID
\item Also very common to use the \alert{standard form} with time constants $\TTi$ and $\TTd$
\item (If you ever see the `series form', that's different again)
\item Two-degree of freedom PID controllers are also possible
\item Cf.\ Matlab's \texttt{pid} / \texttt{pidstd} / \texttt{pid2} / \texttt{pidstd2}, and Simulink PID blocks
\end{itemize}

\end{frame}

\begin{frame}
\frametitle{Generally cannot use a straight derivative}
\[
 u(t) = \kkp e(t) + \kki\int_0^t e(\tau)\dee \tau + \kkd\frac{\dee e(t)}{\dee t}
\]
\[
 \frac{U(s)}{E(s)} = \kkp + \frac{k_i}{s} + \kkd s = \frac{\kkd s^2 + \kkp s + \kki}{s}
\]
Notice anything? Therefore:
\[
   \frac{U_{\text{d}}(s)}{E(s)} = \frac{\kkd s}{1 + \TTf s}
\]
\end{frame}

\SUMMARYFRAME
\FINALE

\end{document}



